\section{Sharp bounds for Boolean phase lines}
\label{bound:section}

The enhanced family also attains the largest radius possible for a broader
class of phase deformations generated by real functions taking at most
two values on the physical input strings; constant generators are included.
Subtracting the generator's value at $z=0$ removes a global phase.
A constant generator then becomes zero; otherwise, rescaling or reversing
$t$ makes the two values $0$ and $1$. We can therefore use a Boolean
generator, without a compression assumption or a condition on the
background state, and write
\begin{equation}
 \begin{gathered}
 M_d(t)=C_n\operatorname{diag}_{z}
 \bigl(i^{f_0(z)}e^{itd(z)}\bigr)H^{\otimes n},\\
 d:\mathbb F_2^n\longrightarrow\{0,1\},\qquad d(0)=0.
 \end{gathered}
 \label{bound:line}
\end{equation}
We call such a line admissible if, on an open interval containing $t=0$,
its local Pauli orbits form regular tetrahedra with a common radius.
This requirement refers to the complete measurement.  The unitary
factorization already guarantees orthonormality for every real $t$.

When the signed local components share a value $m_d(t)$, the common
Bloch radius is
\begin{equation}
 r_d(t)=\sqrt3\,|m_d(t)|.
 \label{bound:radius}
\end{equation}
Intervals with $m_d\ge0$ retain the original tetrahedral orientation.

\begin{theorem}[Four laws and the sharp Boolean-line bound]
\label{bound:theorem}
Let $n\ge3$ and let $d$ be admissible.  Set $c_n=2^{1-n}$.
All signed local components have one common value $m_d(t)$ for every
real $t$, and, for one fixed pair $(s,k)\in\{0,1\}^2$,
\begin{equation}
 m_d(t)=c_n\bigl[(1-s)+s\cos t-k\sin t\bigr].
 \label{bound:four-laws}
\end{equation}
Thus the four possible functional forms of $m_d/c_n$ are
$1$, $\cos t$, $1-\sin t$, and $\cos t-\sin t$.
Not every form occurs in every dimension.
For every $n\ge2$, with $r_E(n)=\sqrt3\,c_n$, the exact maximum over admissible lines is
\begin{equation}
 \max_{\substack{d\ \text{admissible}\\t\in\mathbb R}}r_d(t)
 =\begin{cases}
 r_E(n),&n=2,3,4,\\
 2r_E(n),&n\ge5.
 \end{cases}
 \label{bound:sharp-max}
\end{equation}
Both maxima are attained.
\end{theorem}

At the base point all these components equal $c_n>0$.
Regularity therefore makes them equal on a neighborhood of the base
point.  Each component is a linear combination of $1$, $\cos t$, and
$\sin t$, so equality on that neighborhood implies equality for all
real $t$.

The essential restriction comes from comparing coefficients for the two
classes of pulled-back Pauli operators.  Put $q=d+\delta_0$, so that
$q$ is Boolean and $q(0)=1$, and use the Walsh coefficients and Boolean
derivatives defined in Eq.~\eqref{core:walsh-def}.
For class I, the coefficient of $\cos t$ is $c_nq(a)$; for class II,
the coefficient of $\sin t$ is $-c_nq(a)$.  Since both classes occur
for $n\ge3$, equality of the components forces a common flip value
$s$ in class I and a common flip value $k$ in class II.  The remaining
coefficients agree precisely when
\begin{equation}
 \begin{aligned}
 q(a)&=s,& W_q(h)&=c_n(s+k-1)&&({\rm I}),\\
 q(a)&=k,& W_{D_aq}(h)&=2c_n(s+k-1)&&({\rm II}).
 \end{aligned}
 \label{bound:conditions}
\end{equation}
These conditions give Eq.~\eqref{bound:four-laws} and are also
sufficient.  For the enhanced law $1-\sin t$, they force both Walsh
coefficients to vanish, so Eq.~\eqref{core:boolean-background} makes the
background identically zero.

The maximum absolute values of the four normalized functions are,
respectively, $1$, $1$, $2$, and $\sqrt2$.  This proves the upper bound
$2r_E$ for every $n\ge3$.  The low-dimensional restrictions are stronger.
For $n=3$, two class-I character equations force $s=k=0$.
For $n=4$, four such equations sum to $4k=0$, leaving only the constant
and cosine laws.  The calculation for $n=2$, where class II is absent,
also permits only these two laws.  Appendix~\ref{bound:appendix}
gives the coefficient calculation and these analytic exclusions.
The base point attains the low-dimensional maximum, while the unified
family attains $2r_E$ at $t=-\pi/2$ for every $n\ge5$.

Consequently, $n=5$ is the smallest system size permitting an
increase of the radius along a phase line through the specified base
point with a two-valued real generator. General multivalued phase lines
and curved paths require a separate analysis; the case $n=4$ below
makes this distinction concrete.
